% GENERATED FILE. Edit canonical lessons, then run npm run generate:books.
% canonical-source-hash: fnv1a64:78edaab0b2298346
% canonical-lessons: PT-C278-noticiario, PT-C278-telejornal, PT-C278-apresentador, PT-C278-em-direto, PT-C278-transmissao

\chapter{O noticiário da noite}
\label{ch:pt-o-noticiario-da-noite}
\hlchaptermodality{278}

\begin{chapteropening}
\textbf{By the end of this chapter:} \emph{I can talk about the news on radio and TV: the bulletin, the presenter, live coverage and the broadcast.}

Complete the last lesson of chapter 278: I can talk about the news on radio and TV: the bulletin, the presenter, live coverage and the broadcast.
\end{chapteropening}

\section[o noticiário]{o noticiário — the news bulletin, the news programme}

\label{lesson:PT-C278-noticiario}



\begin{quote}
Before the new one: say the Portuguese for daily, then the Portuguese for beforehand.
\end{quote}

\subsection*{You'll want to know: o noticiário}
\textbf{o noticiário} — \textquotedblleft{}the news bulletin, the news programme\textquotedblright{}.

\begin{quote}
Ouço o noticiário na rádio todas as manhãs.
\end{quote}

\begin{quote}
\emph{I listen to the news on the radio every morning.}
\end{quote}

On radio or TV; a single piece of news is \emph{uma notícia}.

\begin{cousinweb}
\emph{noticiário} is built on \emph{notícia} \textquotedblleft{}news\textquotedblright{}, from Latin \emph{nōtitia} \textquotedblleft{}knowledge, a notice\textquotedblright{}, from \emph{nōscere} \textquotedblleft{}to get to know\textquotedblright{} (Lewis \& Short).
\end{cousinweb}

\begin{grammarlens}[title={What you've built}]
One more word for telling, explaining and describing at B1.
\end{grammarlens}

\subsection*{Your turn}
\begin{itemize}
\raggedright
  \item \emph{Say it:} \emph{o noticiário}
  \item \emph{Say it:} \emph{o noticiário}, once more
  \item \emph{Say it:} a sentence of your own with \emph{o noticiário}
  \item \emph{Recall:} say the Portuguese for daily, then the Portuguese for beforehand, then say \emph{o noticiário} again
\end{itemize}

\begin{tcolorbox}[breakable,colback=teal!4,colframe=teal!35!black,title={Before you move on}]
What does o noticiário mean? (\textquotedblleft{}The news bulletin, the news programme\textquotedblright{}.) Say it once more.
\end{tcolorbox}
\section[o telejornal]{o telejornal — the TV news}

\label{lesson:PT-C278-telejornal}



\begin{quote}
Before the new one: say the Portuguese for beforehand, then the Portuguese for the news bulletin.
\end{quote}

\subsection*{You'll want to know: o telejornal}
\textbf{o telejornal} — \textquotedblleft{}the TV news\textquotedblright{}.

\begin{quote}
O telejornal começa às oito.
\end{quote}

\begin{quote}
\emph{The TV news starts at eight.}
\end{quote}

Built on \emph{televisão} and \emph{jornal}.

\begin{cousinweb}
\emph{telejornal} joins \emph{tele-}, from Greek \emph{tēle} \textquotedblleft{}far off\textquotedblright{}, and \emph{jornal}, from Latin \emph{diurnālis} \textquotedblleft{}daily\textquotedblright{}, built on \emph{diēs} \textquotedblleft{}day\textquotedblright{} (Lewis \& Short).
\end{cousinweb}

\begin{grammarlens}[title={What you've built}]
One more word for telling, explaining and describing at B1.
\end{grammarlens}

\subsection*{Your turn}
\begin{itemize}
\raggedright
  \item \emph{Say it:} \emph{o telejornal}
  \item \emph{Say it:} \emph{o telejornal}, once more
  \item \emph{Say it:} a sentence of your own with \emph{o telejornal}
  \item \emph{Recall:} say the Portuguese for beforehand, then the Portuguese for the news bulletin, then say \emph{o telejornal} again
\end{itemize}

\begin{tcolorbox}[breakable,colback=teal!4,colframe=teal!35!black,title={Before you move on}]
What does o telejornal mean? (\textquotedblleft{}The TV news\textquotedblright{}.) Say it once more.
\end{tcolorbox}
\section[o apresentador]{o apresentador — the presenter, the host}

\label{lesson:PT-C278-apresentador}



\begin{quote}
Before the new one: say the Portuguese for the news bulletin, then the Portuguese for the TV news.
\end{quote}

\subsection*{You'll want to know: o apresentador}
\textbf{o apresentador} — \textquotedblleft{}the presenter, the host\textquotedblright{}.

\begin{quote}
O apresentador do programa é muito simpático.
\end{quote}

\begin{quote}
\emph{The programme’s host is very friendly.}
\end{quote}

A woman is \emph{a apresentadora}.

\begin{cousinweb}
\emph{apresentador} is built on \emph{apresentar} \textquotedblleft{}to present\textquotedblright{}, from Late Latin \emph{praesentāre} \textquotedblleft{}to show, to hold out\textquotedblright{}, from \emph{praesēns} \textquotedblleft{}present\textquotedblright{} (Lewis \& Short).
\end{cousinweb}

\begin{grammarlens}[title={What you've built}]
One more word for telling, explaining and describing at B1.
\end{grammarlens}

\subsection*{Your turn}
\begin{itemize}
\raggedright
  \item \emph{Say it:} \emph{o apresentador}
  \item \emph{Say it:} \emph{o apresentador}, once more
  \item \emph{Say it:} a sentence of your own with \emph{o apresentador}
  \item \emph{Recall:} say the Portuguese for the news bulletin, then the Portuguese for the TV news, then say \emph{o apresentador} again
\end{itemize}

\begin{tcolorbox}[breakable,colback=teal!4,colframe=teal!35!black,title={Before you move on}]
What does o apresentador mean? (\textquotedblleft{}The presenter, the host\textquotedblright{}.) Say it once more.
\end{tcolorbox}
\section[em direto]{em direto — live (on air)}

\label{lesson:PT-C278-em-direto}



\begin{quote}
Before the new one: say the Portuguese for the TV news, then the Portuguese for the presenter.
\end{quote}

\subsection*{You'll want to know: em direto}
\textbf{em direto} — \textquotedblleft{}live (on air)\textquotedblright{}.

\begin{quote}
O jogo passa em direto na televisão.
\end{quote}

\begin{quote}
\emph{The match is on live on television.}
\end{quote}

In Brazil: \emph{ao vivo}. Before the 1990 spelling reform, \emph{em directo}.

\begin{cousinweb}
\emph{direto} is a learned borrowing from Latin \emph{dīrēctus} \textquotedblleft{}straight, direct\textquotedblright{}, the past participle of \emph{dīrigere} \textquotedblleft{}to set straight\textquotedblright{}, built on \emph{regere} \textquotedblleft{}to guide\textquotedblright{} (Lewis \& Short); the inherited form of the same word is \emph{direito}.
\end{cousinweb}

\begin{grammarlens}[title={What you've built}]
One more word for telling, explaining and describing at B1.
\end{grammarlens}

\subsection*{Your turn}
\begin{itemize}
\raggedright
  \item \emph{Say it:} \emph{em direto}
  \item \emph{Say it:} \emph{em direto}, once more
  \item \emph{Say it:} a sentence of your own with \emph{em direto}
  \item \emph{Recall:} say the Portuguese for the TV news, then the Portuguese for the presenter, then say \emph{em direto} again
\end{itemize}

\begin{tcolorbox}[breakable,colback=teal!4,colframe=teal!35!black,title={Before you move on}]
What does em direto mean? (\textquotedblleft{}Live (on air)\textquotedblright{}.) Say it once more.
\end{tcolorbox}
\section[a transmissão]{a transmissão — the broadcast, the transmission}

\label{lesson:PT-C278-transmissao}



\begin{quote}
Before the new one: say the Portuguese for the presenter, then the Portuguese for live.
\end{quote}

\subsection*{You'll want to know: a transmissão}
\textbf{a transmissão} — \textquotedblleft{}the broadcast, the transmission\textquotedblright{}.

\begin{quote}
A transmissão do concerto começou tarde.
\end{quote}

\begin{quote}
\emph{The broadcast of the concert started late.}
\end{quote}

Also passing on an illness: \emph{a transmissão do vírus}.

\begin{cousinweb}
\emph{transmissão} was borrowed from Latin \emph{trānsmissiō} \textquotedblleft{}a sending across\textquotedblright{} (EtymDB), from \emph{trānsmittere}, built on \emph{mittere} \textquotedblleft{}to send\textquotedblright{}.
\end{cousinweb}

\begin{grammarlens}[title={What you've built}]
One more word for telling, explaining and describing at B1.
\end{grammarlens}

\subsection*{Your turn}
\begin{itemize}
\raggedright
  \item \emph{Say it:} \emph{a transmissão}
  \item \emph{Say it:} \emph{a transmissão}, once more
  \item \emph{Say it:} a sentence of your own with \emph{a transmissão}
  \item \emph{Recall:} say the Portuguese for the presenter, then the Portuguese for live, then say \emph{a transmissão} again
\end{itemize}

\begin{tcolorbox}[breakable,colback=teal!4,colframe=teal!35!black,title={Before you move on}]
What does a transmissão mean? (\textquotedblleft{}The broadcast, the transmission\textquotedblright{}.) Say it once more.
\end{tcolorbox}
